\documentclass[10.5pt]{article}
\usepackage[letterpaper,margin=0.76in]{geometry}
\usepackage[T1]{fontenc}
\usepackage{lmodern}
\usepackage{microtype}
\usepackage{amsmath,amssymb,amsthm,mathtools}
\usepackage{enumitem}
\usepackage{xcolor}
\usepackage{titlesec}
\usepackage{hyperref}

\definecolor{ink}{HTML}{1F2933}
\definecolor{muted}{HTML}{66737F}
\definecolor{accent}{HTML}{394B59}
\hypersetup{colorlinks=true,linkcolor=accent,urlcolor=accent,citecolor=accent}
\color{ink}

\titleformat{\section}{\large\bfseries\color{accent}}{}{0pt}{}
\titlespacing*{\section}{0pt}{0.75em}{0.25em}
\setlength{\parindent}{0pt}
\setlength{\parskip}{0.40em}
\setlist{nosep,leftmargin=1.35em}

\newtheorem{proposition}{Proposition}
\newtheorem*{remark}{Remark}

\title{\vspace{-2.3em}\textbf{The Relational Shape of Structural Complexity\\and Neighborhood Density}}
\author{J. R. Landers}
\date{}

\begin{document}
\maketitle
\vspace{-1.25em}

\begin{abstract}
A mathematical object becomes more richly specified as metadata is added to its description. A natural question is whether this additional structure also changes the object's \emph{relational position}: does a more detailed object tend to lie farther from the other objects in a population, and therefore occupy a sparser neighborhood? For a broad class of additive distances the answer is exact rather than heuristic. Additional distinguishing coordinates can only increase pairwise separation, while the density of every fixed-radius neighborhood can only decrease. The effect admits a direct information-theoretic interpretation.
\end{abstract}

\section*{Objects, descriptions, and distance}
Let
\[
\mathcal O=\{O_1,\ldots,O_K\}
\]
be a finite collection of $K$ objects. At descriptive complexity $N$, object $O_i$ is represented by metadata
\[
O_i^{(N)}=(o_{i1},\ldots,o_{iN}),
\]
where $o_{in}$ denotes the value of the $n$th metadata coordinate. Assume separation is measured by an additive nonnegative distance
\[
\delta_N(O_i,O_j)=\sum_{n=1}^{N} d_n(o_{in},o_{jn}),
\qquad d_n(x,y)\ge 0,
\]
with each $d_n$ measuring dissimilarity along coordinate $n$. The next level of structural complexity appends one further coordinate, so that
\[
\delta_{N+1}(O_i,O_j)
=\delta_N(O_i,O_j)+d_{N+1}(o_{i,N+1},o_{j,N+1}).
\]
Thus added structure acts geometrically: it contributes another nonnegative direction along which two objects may differ.

For a fixed object $O_i$, define its mean relational separation
\[
D_N(i)=\frac{1}{K-1}\sum_{j\ne i}\delta_N(O_i,O_j).
\]
Also define its normalized radius-$r$ neighborhood density by
\[
\rho_N(i;r)=\frac{1}{K-1}\sum_{j\ne i}
\mathbf 1\!\left\{\delta_N(O_i,O_j)\le r\right\},
\]
where $\mathbf 1\{\cdot\}$ is the indicator function. The quantity $\rho_N(i;r)$ is simply the fraction of the remaining objects lying within distance $r$ of $O_i$.

\begin{proposition}[Structural separation and neighborhood thinning]
For every $i$, every $j\ne i$, and every radius $r\ge 0$,
\[
\delta_{N+1}(O_i,O_j)\ge \delta_N(O_i,O_j),
\]
and consequently
\[
D_{N+1}(i)\ge D_N(i),
\qquad
\rho_{N+1}(i;r)\le \rho_N(i;r).
\]
If the appended metadata distinguishes $O_i$ from at least one other object, then $D_{N+1}(i)>D_N(i)$ whenever the corresponding coordinate distance is strictly positive.
\end{proposition}

\begin{proof}
Because $d_{N+1}\ge 0$,
\[
\delta_{N+1}(O_i,O_j)-\delta_N(O_i,O_j)
=d_{N+1}(o_{i,N+1},o_{j,N+1})\ge 0.
\]
Averaging this inequality over $j\ne i$ gives $D_{N+1}(i)\ge D_N(i)$. Since no pairwise distance decreases, an object that lies outside a radius-$r$ neighborhood at level $N$ cannot move into that neighborhood merely by appending a nonnegative distance contribution. Hence every indicator in $\rho_N(i;r)$ is nonincreasing, and so is their average. Strict growth of $D_N(i)$ follows if at least one added contribution is positive.
\end{proof}

The proposition gives a simple relational picture: richer descriptions do not merely enlarge the internal representation of an object. When the new structure is distinguishing, they push that object outward relative to the population, thinning its fixed-radius neighborhood.

\section*{The information carried by a new coordinate}
For categorical metadata, the connection to information is especially transparent. Suppose coordinate $n$ takes values in a finite alphabet $\mathcal A_n$ with probabilities $p_n(x)$, and use Hamming distance,
\[
d_n(x,y)=\mathbf 1\{x\ne y\}.
\]
For two independently sampled objects,
\[
\Pr(o_{in}\ne o_{jn})
=1-\sum_{x\in\mathcal A_n}p_n(x)^2.
\]
Therefore
\[
\mathbb E\,\delta_N(O_i,O_j)
=\sum_{n=1}^{N}\left(1-\sum_{x\in\mathcal A_n}p_n(x)^2\right).
\]
Each coordinate contributes exactly its probability of distinguishing two sampled objects. A constant coordinate contributes zero; a highly diverse coordinate contributes more. The term $\sum_x p_n(x)^2$ is the collision probability and is directly related to R\'enyi entropy of order $2$,
\[
H_2(p_n)=-\log\!\left(\sum_x p_n(x)^2\right).
\]
Thus the growth of expected distance is governed not by metadata count alone, but by the \emph{distinguishing information} carried by the appended structure.

In the uniform $q$-ary case, $p_n(x)=1/q$, so
\[
\mathbb E\,\delta_N=N\left(1-\frac1q\right).
\]
Expected separation grows linearly with description length. For real-valued independent coordinates of variance $\sigma^2$, squared Euclidean distance yields the analogous identity
\[
\mathbb E\,\delta_N^2=2N\sigma^2,
\]
so the typical Euclidean scale grows on the order of $\sqrt N$.

\section*{What the principle does -- and does not -- say}
The essential distinction is between \emph{more metadata} and \emph{more distinguishing structure}. Appending the same value to every object increases description length but leaves all distances unchanged. Redundant coordinates may likewise add little genuinely new separation. Moreover, dimension-normalized distances can behave differently because the normalization itself changes with $N$.

Under a distance that accumulates nonnegative coordinate-wise differences, however, the relational consequence is sharp:
\[
\boxed{
\text{distinguishing structural complexity increases separation and thins local neighborhoods.}
}
\]
This suggests viewing structural complexity not only as an intrinsic property of an object, but as a geometric property of its position among alternatives. As an object becomes more specifically described, its identity is expressed partly through the progressive disappearance of nearby peers.

\end{document}
